\documentclass[12pt]{article}
\usepackage{zed-csp,graphicx,color}%from
\begin{document}

\begin{titlepage}
 \begin{figure}[h]
  \centerline{\small MAKERERE 
  \includegraphics[width=0.1\textwidth]{muk_log} UNIVERSITY}
\end{figure}
\centerline{COLLEGE OF COMPUTING AND INFORMATIC SCIENCES}
\paragraph{•}
\centerline{DEPARTMENT OF COMPUTER SCIENCE\\}
\paragraph{•}

\centerline{COURSEWORK: RESEARCH METHODOLOGY(BIT 2207)\\}
\paragraph{•}

\centerline{LECTURER: MR.ERNEST MWEBAZE}
\paragraph{•}

\centerline{TOPIC\\}HOW THE RISE OF COHABITATION CAN BE PUT TO HALT AROUND KAMPALA SUBURBS. \\
\paragraph{•}
\centerline{COMPILED BY: \
 NAKUBULWA AGNES}
 \paragraph{•}
\centerline{STUDENT NUMBER : 216017951}
\paragraph{•}
\centerline{REGISTRATION NUMBER:16/U/8895/PS}
\paragraph{•}
\begin{flushright}
    Signature ....................\\
    DATE: $ FEBRUARY,8^{TH},2018$
\end{flushright}

\end{titlepage}
\pagenumbering{roman}
\tableofcontents
\newpage
\pagenumbering{arabic}
\section{Introduction} 
Cohabitation is the state of living together and having a sexual relationship without being married.
Cohabiting in Uganda is on remarkably a high rise each passing time. Kampala the capital city in Uganda has one of the greatest cohabiting couple's populations.
This report serve to examine the causes of the rise, possible effects and dangers the arise with cohabiting. Also the possible actions that can be undertaken to reduce or eliminate cohabiting.

\section{Scope of the research}
The research was generally carried out in Uganda, within the suburbs of Kampala. Within the period 2017 to date.
\section{Experimental method}
The information for this report was sourced from different human rights and early marriage articles, Newspapers, individuals plus the internet.
\section{Results}
The following findings were made about the subject matter. The reasons among which include the following:-

\subsection{Need for assistance and attention}
Women and girls especially, often think that the closer one lives with their supposed spouse, the better they are cared for as well as their need being attended too quickly.
\subsection{Pressure}
This resulting from peers leading one to live with a spouse because their friends do too as well. Also some parents put in place a particular age at which their children must have left home.
\subsection{Fear of life commitment}
Many youth are afraid of vowing life commitments. They are not sure if they are with the right partner meanwhile as told by Justine Nakyanzi a 21 year old Nsambya resident.

\subsection{poverty}
	
This pushing many into temporary quick commitment. Most cohabit with the hope of rescue from neediness and vulnerability by their rich partners who in the long run are never interested in commitment.

\section{Suggestion Measures}
Man and woman are meant to be happy together within an official bond of marriage witnessed by a body of the elders and friends or relatives. I suggest the youth do not just give in to marriage none official, also the elders and parents should do their best to advise children right from adolescent time. The government can enforce strict laws on couples living together without the bond of official marriage.

\section{Conclusion}
Cohabiting comes along with much more negativity than positivity. For most couples, their supposed spouses are never known to their families hence no one to freely consult in case anything went wrong, less support or never at all from either side's family members as people are not officially known as couple.






\end{document}